Machine learning models are commonly used for predictions that inform downstream decision-making.
  Recent benchmarks of end-to-end training have broadly focused on methods that minimize tractable surrogates of an ideal yet intractable loss. Less attention has been paid to perturbation methods that average over noisy inputs to obtain informative gradients for the ideal loss.
  In this paper, we make the case that perturbation methods can sometimes match or outperform widely used surrogate approaches in several decision-focused learning benchmarks, even if the chosen prediction model is misspecified.
  However, these gains require careful tuning of the noise-level hyperparameter and do not translate to all tasks. We conclude with realistic guidance for practitioners to get the most out of these methods.

%  Recent research has demonstrated that training models end-to-end on the prediction and decision task leads to superior outcomes. Recently, a family of methods has been created to solve these problems, each with their own advantages, guarantees, and trade-offs. However, there exists no comprehensive benchmark or guidelines for method selection for the machine learning practitioner. Here we present such a benchmark, with heuristics for hyperparameter selection and guidelines for method choice. Additionally, we seek to explain the observed success of some models via the Gradient Signal to Noise Ratio (GSNR).
